\chapter{Enunciado del problema}

Se plantea desarrollar una descripción comparativa de los protocolos de intercambio de claves de variable discreta. Para ello se requiere seleccionar un marco común de análisis y aplicarlo a todos los protocolos. Dicho marco debe recoger las características principales, fundamentos, puntos fuertes, puntos débiles y una explicación práctica con ejemplos.

En este Trabajo se interpreta el problema como una comparación razonada de protocolos DV-QKD. No se pretende demostrar formalmente la seguridad composable de cada esquema, sino explicar sus mecanismos esenciales y su posicionamiento. Por ello se combinan referencias fundacionales con una lectura práctica orientada a la toma de decisiones: qué aporta cada protocolo, qué supuestos exige, qué ataques intenta mitigar y en qué contexto puede ser preferible.